\cleardoublepage
\section{Ratios}

Welcome to Unit 3! In this section, we'll learn how to compare amounts using ratios. A \textbf{ratio} is a comparison between two numbers or amounts. It tells us how much of one thing there is compared to another. Ratios help us describe relationships and patterns, and they appear in recipes, sports, and even classrooms.

**Example:**  
If there are 2 apples and 3 oranges, the ratio of apples to oranges is **2 to 3**.

\subsection*{Ways to Write Ratios}

There are several ways to write ratios:

\begin{itemize}
  \item With a colon: \textbf{2:3}
  \item As a fraction: \textbf{2/3}
  \item In words: \textbf{2 to 3}
\end{itemize}

All three ways mean the same thing.

\subsection*{Part-to-Part and Part-to-Whole}

\begin{itemize}
  \item \textbf{Part-to-Part:} Compares two individual parts (e.g., apples to oranges → 2:3).
  \item \textbf{Part-to-Whole:} Compares one part to the total (e.g., apples to total fruits → 2:5).
\end{itemize}

**Example:**  
If a class has 12 boys and 8 girls:
\begin{itemize}
  \item The ratio of boys to girls is \textbf{12:8} → simplified to \textbf{3:2}.
  \item The ratio of girls to total students (20) is \textbf{8:20} → simplified to \textbf{2:5}.
\end{itemize}

Ratios are used to mix paint colors, create scale models, adjust ingredients in recipes, and describe odds in games!


\cleardoublepage
\subsection*{Short Question (English)}
There are 6 red marbles and 4 blue marbles in a bag. What is the ratio of red marbles to blue marbles?

\fullpageimage{images/q3_1_1-eng.jpg}

\newpage
\subsection*{Short Question (Korean)}
There are 6 red marbles and 4 blue marbles in a bag. What is the ratio of red marbles to blue marbles?

\fullpageimage{images/q3_1_1-kor.jpg}

\newpage
\subsection*{Short Question (Answer)}
There are 6 red marbles and 4 blue marbles in a bag. What is the ratio of red marbles to blue marbles?

\textbf{Answer:} \textbf{6:4}, which can be simplified to \textbf{3:2}.

\cleardoublepage
\subsection{title_here (English)}

question here

\fullpageimage{images/q3_1_2-eng.jpg}


\newpage
\subsection{title_here (Korean)}

question here

\fullpageimage{images/q3_1_2-kor.jpg}


\newpage
\subsection{title_here (Answer)}

question here

\cleardoublepage
\subsection{Challenge Question (English)}
A recipe for bibimbap calls for 2 cups of rice and 3 cups of vegetables. If you want to make enough for twice as many people, keeping the same ratio, how many cups of rice and vegetables do you need?

\fullpageimage{images/q3_1_3-eng.jpg}

\newpage
\subsection{Challenge Question (Korean)}
A recipe for bibimbap calls for 2 cups of rice and 3 cups of vegetables. If you want to make enough for twice as many people, keeping the same ratio, how many cups of rice and vegetables do you need?

\fullpageimage{images/q3_1_3-kor.jpg}

\newpage
\subsection{Challenge Question (Answer)}
A recipe for bibimbap calls for 2 cups of rice and 3 cups of vegetables. If you want to make enough for twice as many people, keeping the same ratio, how many cups of rice and vegetables do you need?

\textbf{Answer:}\\
Rice: $2 \times 2 = 4$ cups\\
Vegetables: $3 \times 2 = 6$ cups

